\section{A shorter slope comparison for a family of integral maps}
\label{cs-section}

The final arithmetic comparison in the proposed argument\cite{alpo-central-comparison} has two distinct
parts. One constructs integral maps, controls their heights and local
divisibility, and proves that their complex realizations are nearly isotropic.
The other turns these facts into an inequality between slopes. We give a
short proof of the second part. Its hypotheses are concrete; establishing
them for the required automorphic family remains a separate problem.

\subsection{Slopes and the height of a map}

Let $K$ be a number field of degree $d$. A Hermitian $\mathcal O_K$-lattice
$E$ is a finite projective $\mathcal O_K$-module with a positive definite
Hermitian norm at each complex embedding, compatible with conjugation.
Its rank is its dimension over $K$. Give its real Minkowski space the
metric in which a complex place contributes twice the real part of the
Hermitian form. Write $\operatorname{covol}(E)$ for the resulting real
covolume and put
\[
 \widehat{\deg}(E)
 =\frac1d\left(-\log\operatorname{covol}(E)
                  +\frac{\operatorname{rk}E}{2}\log|\Delta_K|\right),
 \qquad \mu(E)=\frac{\widehat{\deg}(E)}{\operatorname{rk}E}.
\]
Submodules are given the restricted metrics and quotients the quotient
metrics. A submodule is called saturated if its quotient is torsion free.
For a nonzero integral linear map $f:E\to F$, its height is
\[
 h(f)=\frac1d\sum_{v\mid\infty}n_v\log\|f_v\|_{\rm op},
 \qquad n_v=[K_v:\mathbb R].
\]
We allow the zero map and assign it height $-\infty$.

We recall the slope facts that will be used. Every nonzero $E$ has a
filtration by saturated submodules
\[
 0=E_0\subset E_1\subset\cdots\subset E_a=E
\]
whose successive quotients are semistable and have strictly decreasing
slopes. Semistable means that no nonzero saturated submodule has larger
slope than the whole lattice. List the successive slopes with their
rank multiplicities as
$\lambda_1(E)\ge\cdots\ge\lambda_r(E)$.
Then their mean is $\mu(E)$; write
$\mu_{\max}(E)=\lambda_1(E)$ and
$\mu_{\min}(E)=\lambda_r(E)$. For a real $t$, let $E_{\ge t}$ denote the
step consisting of all successive quotients with slope at least $t$.

For completeness, the construction uses only covolumes. In each fixed
rank, the degrees of saturated submodules are bounded above and attain
their maximum: their underlying primitive integral sublattices have
nonzero exterior products in a fixed discrete lattice, and there are
only finitely many such products of bounded norm. Choose a submodule
of maximal slope and, among those, maximal rank. It is unique. Indeed,
for two saturated submodules $U,V$ the determinant formula for an exact
sequence and the decrease of quotient norms give
\[
 \widehat{\deg}U+\widehat{\deg}V
 \le\widehat{\deg}(U\cap V)
       +\widehat{\deg}\bigl((U+V)^{\rm sat}\bigr).
\]
If $U,V$ have maximal slope, both terms on the right are at most that
slope times their respective ranks; the ranks add to
$\operatorname{rk}U+\operatorname{rk}V$. Thus their saturated sum also has
maximal slope. Maximality of the chosen rank proves uniqueness.
Repeat in the quotient. The slopes strictly decrease, since otherwise
the inverse image of the next maximal submodule contradicts the choice
at the preceding step. The process terminates after at most $r$ steps.
Degree additivity proves the assertion about the mean. It also shows
that every nonzero torsion-free quotient, with its quotient metrics,
has slope at least $\mu_{\min}(E)$ and every
nonzero submodule has slope at most $\mu_{\max}(E)$: intersect a submodule,
or filter a quotient, by the displayed filtration, and use semistability
on its successive images. Saturation and quotient norms can only improve
the needed inequalities.

\begin{lemma}[Slopes under an integral map]\label{cs-functoriality}
If $f:E\to F$ is nonzero, then
\[
 \mu_{\min}(E)\le\mu_{\max}(F)+h(f).
\]
If $h(f)\le0$, then $f(E_{\ge t})\subset F_{\ge t}$ for every real $t$.
\end{lemma}
\begin{proof}
Let $Q=E/\ker f$ and let $J$ be the saturation of $f(E)$ in $F$.
Put $q=\operatorname{rk}Q=\operatorname{rk}J$. On $q$th exterior powers
the operator norm is at most $\|f_v\|_{\rm op}^{q}$.
The image lattice has finite index in $J$; its index contributes a
nonnegative logarithm to its covolume. Multiplying the determinant
inequalities at the infinite places therefore gives
\[
 \widehat{\deg}(Q)\le\widehat{\deg}(J)+q h(f).
\]
Divide by $q$ and use the preceding submodule and quotient properties.
For the second assertion, compose the restriction to $E_{\ge t}$ with
$F\to F/F_{\ge t}$. Restriction and orthogonal quotient do not increase
operator norms. If this composite were nonzero, the first assertion
would give $t\le\mu_{\min}(E_{\ge t})
\le\mu_{\max}(F/F_{\ge t})<t$, a contradiction.
\end{proof}

\subsection{The matrix estimate}

For a complex matrix $A$, put
$\|A\|_{\mathrm{HS}}^2=\operatorname{tr}(A^*A)$.
An orthogonal projection always refers to the ordinary Euclidean inner
product, independently of the arithmetic metrics above.

\begin{lemma}[Images of subspaces]\label{cs-dimension}
Let $A_i:\mathbb C^{d_i}\to\mathbb C^D$ be a finite family of matrices,
where $D\ge1$ and $R=\sum_i d_i>0$.
Set $Q=\sum_i A_iA_i^*$ and suppose that
\begin{align}
 \sum_i\|A_i^*A_i-I\|_{\mathrm{HS}}^2&\le eR,
       &0\le e\le\tfrac14,\label{cs-diagonal}\\
 \operatorname{tr}(Q^2)&\le
       \frac{1+\eta}{D}(\operatorname{tr}Q)^2,
       &\eta\ge0.\label{cs-isotropic}
\end{align}
If $V_i\subset\mathbb C^{d_i}$ and all $A_i(V_i)$ lie in a subspace
$W\subset\mathbb C^D$, put $t=\sum_i\dim V_i$ and $r=\dim W$. Then
\begin{equation}\label{cs-rank-bound}
 \frac rD\ge
 \frac{(t/R-\sqrt e)_+^2}{(1+\eta)(1+\sqrt e)^2},
 \qquad z_+=\max\{z,0\}.
\end{equation}
\end{lemma}
\begin{proof}
Let $P_i$ project onto $V_i$ and put $C=\sum_i A_iP_iA_i^*$.
Then $0\le C\le Q$ and $\operatorname{rank}C\le r$. Cauchy--Schwarz gives
\[
 |\operatorname{tr}Q-R|\le\sqrt{R}
     \left(\sum_i\|A_i^*A_i-I\|_{\mathrm{HS}}^2\right)^{1/2}
 \le\sqrt e\,R
\]
and, using $\sum_i\|P_i\|_{\mathrm{HS}}^2=t$,
\[
 \operatorname{tr}C
 =t+\sum_i\operatorname{tr}\bigl(P_i(A_i^*A_i-I)\bigr)
 \ge t-\sqrt{eRt}\ge t-\sqrt e\,R.
\]
The eigenvalues of $C$ are nonnegative. Thus
$(\operatorname{tr}C)^2\le r\operatorname{tr}(C^2)$.
Moreover
$\operatorname{tr}(Q^2-C^2)
=\operatorname{tr}((Q-C)(Q+C))\ge0$, since the trace of a product
of two positive semidefinite matrices is nonnegative. Combine these
facts with \eqref{cs-isotropic} and
$\operatorname{tr}Q\le(1+\sqrt e)R$. If $t/R\le\sqrt e$, the desired
inequality is simply $r\ge0$; otherwise division gives
\eqref{cs-rank-bound}.
\end{proof}

\begin{proposition}[Slope comparison]\label{cs-slope}
Let $I$ be a finite nonempty index set. Let $E_i$ for $i\in I$, and $F$,
be Hermitian $\mathcal O_K$-lattices with positive ranks
$d_i,n$, put $R=\sum_{i\in I}d_i$, and let $f_i:E_i\to F$ be integral linear maps with
$h(f_i)\le h_i\in\mathbb R$. Suppose that
\begin{equation}\label{cs-numerical}
 \frac1R\sum_i d_i\bigl(\mu(E_i)-h_i\bigr)\ge s,
 \quad \mu_{\max}(E_i)-h_i\le s+u,
 \quad \mu_{\min}(F)\ge s-H,
 \qquad u,H\ge0.
\end{equation}
Fix one complex embedding of $K$. Suppose that $D\ge n$,
$n\ge(1-\delta)D$, and $0\le\delta<1$. Suppose that matrices
$A_i:\mathbb C^{d_i}\to\mathbb C^D$ satisfy
\eqref{cs-diagonal}--\eqref{cs-isotropic}, and that their first $n$ rows
represent the maps $f_i$ in some complex bases of $E_i$ and a single
common basis of $F$. Positive scalar factors depending on $i$ are allowed.
Then
\begin{equation}\label{cs-slope-bound}
 s-\mu(F)\le\frac{2u\log(1+H/u)+(\eta+4\sqrt e)H}{1-\delta}.
\end{equation}
At $u=0$ the expression $u\log(1+H/u)$ means its continuous value $0$.
\end{proposition}
\begin{proof}
Multiply every norm on $E_i$ by $\exp(h_i)$. All of its slopes decrease
by $h_i$, and the height of $f_i$ decreases by $h_i$. Thus the maps now
have nonpositive height. By Lemma~\ref{cs-functoriality}, the source
subspaces with slopes at least $s-z$ map into $F_{\ge s-z}$.

For $z>0$, let $t(z)$ be the sum of the ranks of these source subspaces.
Every adjusted source slope is at most $s+u$, and their weighted mean is
at least $s$. The total of their nonnegative deficits from $s+u$ is at
most $uR$. Each slope smaller than $s-z$ has deficit larger than $u+z$.
Consequently
\[
 a(z):=t(z)/R\ge\frac z{u+z}.
\]
Write $r(z)=\operatorname{rk}F_{\ge s-z}$. The images in $\mathbb C^D$
lie in the direct sum of this $r(z)$-dimensional target subspace and the
last $D-n$ coordinate directions. Lemma~\ref{cs-dimension} gives
\[
 1-\frac{r(z)}n\le\frac1{1-\delta}
 \left(1-\frac{(a(z)-\sqrt e)_+^2}
 {(1+\eta)(1+\sqrt e)^2}\right).
\]
For $0\le a\le1$, $q\ge0$, and $\eta\ge0$ one has
\[
 1-\frac{(a-q)_+^2}{(1+\eta)(1+q)^2}
 \le\eta+4q+2(1-a).
\]
To see this, put $b=(a-q)_+/(1+q)$. Then $0\le b\le a\le1$,
$a-b\le2q$, and
$1-b^2/(1+\eta)\le\eta+(1-a^2)+(a^2-b^2)
\le\eta+2(1-a)+4q$.
It follows that
\[
 1-r(z)/n\le\frac{\eta+4\sqrt e+2u/(u+z)}{1-\delta}.
\]
Every target slope is at least $s-H$. If these slopes are $\nu_j$,
integration of the last inequality therefore yields
\begin{align*}
 s-\mu(F)&\le\frac1n\sum_j(s-\nu_j)_+
   =\int_0^H\left(1-\frac{r(z)}n\right)\,dz\\
 &\le\frac{(\eta+4\sqrt e)H+2u\log(1+H/u)}{1-\delta}.
\end{align*}
When $u=0$, the source dimension is $R$ for every $z>0$ and the same
proof directly gives the stated limiting formula.
\end{proof}

\begin{lemma}[Pairwise estimates suffice]\label{cs-pairwise}
Let $A_i:\mathbb C^{d_i}\to\mathbb C^D$, $D\ge1$, be a finite family,
$R=\sum_i d_i>0$, and $Q=\sum_iA_iA_i^*$.
Suppose $D_*>0$, $\eta_0\ge-1$, $\eta_1\ge0$, $0\le e\le1/4$, and
\[
 \|A_i^*A_j\|_{\mathrm{HS}}^2
       \le(1+\eta_0)d_id_j/D_*\quad(i\ne j),\qquad
 \|A_i^*A_i-I\|_{\mathrm{HS}}^2\le e d_i,
 \qquad D\le(1+\eta_1)D_*.
\]
Then \eqref{cs-diagonal} holds and
\[
 \frac{D\operatorname{tr}(Q^2)}{(\operatorname{tr}Q)^2}
 \le
 \frac{(1+\eta_1)
   \bigl(1+\eta_0+(D_*/R)(1+\sqrt e)^2\bigr)}{(1-\sqrt e)^2}.
\]
Thus any nonnegative $\eta$ at least the right side minus $1$ is
permitted in Proposition~\ref{cs-slope}.
\end{lemma}
\begin{proof}
Cyclicity of trace gives
$\operatorname{tr}(Q^2)=\sum_{i,j}\|A_i^*A_j\|_{\mathrm{HS}}^2$.
The off-diagonal terms sum to at most $(1+\eta_0)R^2/D_*$.
The triangle inequality gives
$\|A_i^*A_i\|_{\mathrm{HS}}\le(1+\sqrt e)\sqrt{d_i}$.
Therefore
\[
 \operatorname{tr}(Q^2)
 \le(1+\eta_0)R^2/D_*+(1+\sqrt e)^2R.
\]
The proof of Lemma~\ref{cs-dimension} gives
$\operatorname{tr}Q\ge(1-\sqrt e)R>0$. Division proves the claim.
\end{proof}

\begin{lemma}[Restriction to subspaces]\label{cs-compression}
Let $B_i:\mathbb C^{r_i}\to\mathbb C^D$, $D\ge1$, be a finite family satisfying the pairwise and
diagonal inequalities of Lemma~\ref{cs-pairwise}, with parameters
$D_*,\eta_0,\eta_1,e_0$. Let
$V_i:\mathbb C^{d_i}\to\mathbb C^{r_i}$ satisfy $V_i^*V_i=I$,
and put $A_i=B_iV_i$, $R_f=\sum_i r_i>0$, $R=\sum_i d_i$.
Suppose $R_f-R\le\theta R_f$, where $0\le\theta<1$, and set
\[
 q=\sqrt{e_0},\qquad
 \kappa=1-\theta-q>0.
\]
Then $R>0$, and the compressed matrices satisfy
\[
 \sum_i\|A_i^*A_i-I\|_{\mathrm{HS}}^2
       \le\frac{e_0}{1-\theta}R,
\]
and, for $Q=\sum_iA_iA_i^*$,
\[
 \frac{D\operatorname{tr}(Q^2)}{(\operatorname{tr}Q)^2}
 \le\frac{(1+\eta_1)
 [1+\eta_0+(D_*/R_f)(1+q)^2]}{\kappa^2}.
\]
After omitting zero-dimensional domains, these are the matrix hypotheses
of Proposition~\ref{cs-slope} whenever
$e_0/(1-\theta)\le1/4$, with the displayed bound determining $1+\eta$.
\end{lemma}
\begin{proof}
Orthogonal compression decreases the Hilbert--Schmidt norm, so
\[
 \|A_i^*A_i-I\|_{\mathrm{HS}}
 =\|V_i^*(B_i^*B_i-I)V_i\|_{\mathrm{HS}}
 \le\|B_i^*B_i-I\|_{\mathrm{HS}}.
\]
This follows also by extending the columns of $V_i$ to an orthonormal
basis and summing the squares of matrix entries. Summing and using
$R\ge(1-\theta)R_f>0$ proves the first assertion.

Put $Q_f=\sum_iB_iB_i^*$ and let $P_i=V_iV_i^*$.
Then $0\le Q\le Q_f$, and $\sum_i\operatorname{rank}P_i=R$.
Cauchy--Schwarz gives
\begin{align*}
 \operatorname{tr}Q
 &=R+\sum_i\operatorname{tr}\bigl(P_i(B_i^*B_i-I)\bigr)\\
 &\ge R-\sqrt{e_0R_fR}\ge(1-\theta-q)R_f.
\end{align*}
Thus $\operatorname{tr}Q\ge\kappa R_f>0$.
As in Lemma~\ref{cs-dimension},
$\operatorname{tr}(Q^2)\le\operatorname{tr}(Q_f^2)$.
Use the upper bound for this last trace from the proof of
Lemma~\ref{cs-pairwise} and divide by $\kappa^2R_f^2$.
\end{proof}

In the source application, $D_*/R_f\le\sigma$,
$\eta_0,\eta_1,\sqrt{e_0}\le\sigma/2$, $0<\sigma\le1/20$,
and $0\le\theta\le1/4$; a negative $\eta_1$ can harmlessly be replaced
by $0$. These imply $\kappa>0$ and $e_0/(1-\theta)<1/4$.
In fact the displayed trace-ratio bound is at most $1+7\sigma+4\theta$.
Indeed, its numerator is at most $1+9\sigma/4$, and
$\kappa\ge1-t$ for $t=\theta+\sigma/2\le11/40$.
For $0\le t\le1/3$ the inequality $(1-t)^{-2}\le1+4t$ follows by
multiplication and $2-7t+4t^2\ge0$. Hence the trace ratio is at most
$(1+9\sigma/4)(1+4\theta+2\sigma)
\le1+7\sigma+4\theta$, using $\theta\le1/4$ and $\sigma\le1/20$.
Thus the passage from the pairwise estimates
on the full matrices to the required estimate on their restrictions
is justified by ordinary matrix inequalities. The quantitative
pairwise estimates themselves are still an input about the theta maps.

% Optional refinement of the independent comparison in central_step.tex.
% This subsection uses the notation and labels introduced there.
\subsection{A refinement using the complementary subspaces}
\label{qrh-sharper-slope-section}

The same hypotheses give a second bound that can be useful when the
source slope spread is much smaller than the permitted target deficit.
The simple estimate in Proposition~\ref{cs-slope} remains available;
one may use whichever of the two bounds is smaller.

\begin{proposition}[Refined slope comparison]
\label{qrh-sharper-slope}
Keep the lattice, map, rank, and numerical hypotheses of
Proposition~\ref{cs-slope}. In its matrix hypotheses, assume
$0\le e<1$ in place of $e\le1/4$, and retain
\eqref{cs-isotropic} with $\eta\ge0$.
Put $q=\sqrt e$ and, for $u,H\ge0$, define
\[
 \mathcal A=
 \frac{u\log(1+H/u)
       +2q\sqrt u\bigl(\sqrt{u+H}-\sqrt u\bigr)}{1-q}.
\]
At $u=0$ this expression means its continuous value $0$;
at $H=0$ it also equals $0$. Then
\begin{equation}\label{qrh-sharper-slope-bound}
 s-\mu(F)\le
 \frac{\mathcal A+\eta H/2+
       \sqrt{\mathcal A\eta H+(\eta H)^2/4}}{1-\delta}.
\end{equation}
If $u=0$, the stronger bound
\[
 s-\mu(F)\le\frac{\eta H}{(1+\eta)(1-\delta)}
\]
holds. In particular, when $e=\eta=0$,
\begin{equation}\label{qrh-exact-slope-bound}
 s-\mu(F)\le\frac{u\log(1+H/u)}{1-\delta}.
\end{equation}
All assertions include $H=0$ and $u=0$ with the stated conventions.
\end{proposition}

\begin{proof}
As in Proposition~\ref{cs-slope}, rescale the source norms by
$\exp(h_i)$ so that all map heights are nonpositive. For $z>0$,
let $V_i(z)$ be the complex realization, in the chosen source basis,
of the subspace of adjusted slopes at least $s-z$. Set
\[
 t(z)=\sum_i\dim V_i(z),\qquad c(z)=1-t(z)/R.
\]
The source slope assumptions give
\begin{equation}\label{qrh-source-complement}
 0\le c(z)\le\frac{u}{u+z}.
\end{equation}
When $u=0$, every adjusted source slope equals $s$, so $c(z)=0$.
By Lemma~\ref{cs-functoriality}, the images of $V_i(z)$ under $A_i$
lie in the subspace $W(z)$ obtained by adjoining the last $D-n$
coordinate directions to $F_{\ge s-z}$. Write
\[
 r(z)=\operatorname{rk}F_{\ge s-z},\qquad
 x(z)=\frac{D-\dim W(z)}D=\frac{n-r(z)}D.
\]

Fix $z$. Let $P_i$ be the orthogonal projection onto $V_i(z)$,
let $P$ be the orthogonal projection onto $W(z)^\perp$, and put
\[
 Q=\sum_i A_iA_i^*,\qquad
 C=\sum_i A_iP_iA_i^*,\qquad T=\operatorname{tr}Q.
\]
Thus $0\le C\le Q$ and $PC=0$. Since $0\le P\le I$, one has
$\operatorname{tr}(PQ)\le\operatorname{tr}(Q-C)$.
Using Cauchy--Schwarz for the sum of the Hilbert--Schmidt inner
products gives
\begin{align*}
 \operatorname{tr}(Q-C)
 &=R-t+\sum_i\operatorname{tr}
       \bigl((I-P_i)(A_i^*A_i-I)\bigr)\\
 &\le R-t+\sqrt{(R-t)eR}
   =R(c+q\sqrt c).
\end{align*}
The diagonal hypothesis also gives $T\ge(1-q)R>0$.
Consequently
\begin{equation}\label{qrh-complement-trace}
 \frac{\operatorname{tr}(PQ)}T
       \le d(z):=\frac{c(z)+q\sqrt{c(z)}}{1-q}.
\end{equation}

The isotropy hypothesis is exactly the variance bound
\[
 \left\|Q-\frac TD I\right\|_{\mathrm{HS}}^2
       \le\frac{\eta T^2}D.
\]
Since $\operatorname{tr}P=Dx$ and
$\|P-xI\|_{\mathrm{HS}}^2=Dx(1-x)$, another application of
Cauchy--Schwarz yields
\[
 \frac{\operatorname{tr}(PQ)}T
 \ge x-\sqrt{\eta x(1-x)}.
\]
Combining this with \eqref{qrh-complement-trace}, and using
$0\le x\le1$, gives
\begin{equation}\label{qrh-complement-dimension}
 x(z)\le d(z)+\sqrt{\eta x(z)}.
\end{equation}

Integrate over $0<z<H$. By \eqref{qrh-source-complement},
\[
 \int_0^H d(z)\,dz\le\mathcal A,
\]
because the integrals of $u/(u+z)$ and
$\sqrt{u/(u+z)}$ are respectively
$u\log(1+H/u)$ and
$2\sqrt u(\sqrt{u+H}-\sqrt u)$.
If $X=\int_0^H x(z)\,dz$, Cauchy--Schwarz in
\eqref{qrh-complement-dimension} therefore gives
\[
 X\le\mathcal A+\sqrt{\eta HX}.
\]
Solving the quadratic inequality for $\sqrt X$ proves
\[
 X\le\mathcal A+\frac{\eta H}{2}
              +\sqrt{\mathcal A\eta H+\frac{(\eta H)^2}{4}}.
\]
Every target slope is at least $s-H$. Hence, exactly as in the
proof of Proposition~\ref{cs-slope},
\[
 s-\mu(F)\le\int_0^H\frac{n-r(z)}n\,dz
       =\frac Dn X\le\frac X{1-\delta},
\]
which proves \eqref{qrh-sharper-slope-bound}.

If $u=0$, then $c(z)=0$ and
$\operatorname{tr}(PQ)=0$. The sharper variance inequality gives
$x\le\sqrt{\eta x(1-x)}$. For $x>0$, squaring and dividing by
$x$ gives $x\le\eta/(1+\eta)$; the same bound holds when $x=0$.
Integrating proves the asserted endpoint improvement. If $e=\eta=0$,
then $\mathcal A=u\log(1+H/u)$ and the general bound is
\eqref{qrh-exact-slope-bound}. When $H=0$, the target slope
assumption already gives $\mu(F)\ge s$, so all displayed bounds
hold without an integration interval.
\end{proof}

The special case $e=\eta=0$ means that the matrices are isometries
on their source spaces and that $Q=(R/D)I$. It gives coefficient
$1$ in front of $u\log(1+H/u)$, in place of the coefficient $2$
in the simpler comparison. For nonzero diagonal error,
\[
 \mathcal A\le
 \frac{u\log(1+H/u)+2\sqrt e\,\sqrt{uH}}{1-\sqrt e},
\]
so this form can reduce the effect of the diagonal error when
$u$ is much smaller than $H$. The mixed term involving $\eta$
is explicit in \eqref{qrh-sharper-slope-bound}; no assertion of
uniform improvement over every parameter choice is needed.
This refinement concerns only the comparison of slopes. It does
not establish the integral automorphic maps, their height bounds,
or their matrix estimates.


\subsection{Discarding an orthogonal subspace}

An exceptional subspace affects a covolume through its degree and the
index of an integral splitting. The following calculation keeps these
two losses explicit.

Let $L$ be a nonzero Hermitian $\mathcal O_K$-lattice and let $P$ be a
$K$-linear idempotent on $L_K$, orthogonal at every infinite place.
Suppose that a positive integer $N$ satisfies $NP(L)\subset L$. Put
\[
 G=L\cap\operatorname{im}P,\qquad B=L\cap\ker P,\qquad
 g=\operatorname{rk}G,\quad b=\operatorname{rk}B,\quad r=g+b,
\]
and assume $g>0$. Let $I=[L:G\oplus B]$ be the finite index.

\begin{proposition}[The exact loss in an orthogonal splitting]
\label{cs-orthogonal-cut}
With $d=[K:\mathbb Q]$, one has
\[
 1\le I\le N^{d\min(g,b)},\qquad
 g\mu(G)=r\mu(L)-b\mu(B)-\frac{\log I}{d}.
\]
The term $b\mu(B)$ is interpreted as zero when $B=0$.
\end{proposition}

\begin{proof}
The decomposition $L_K=\operatorname{im}P\oplus\ker P$ gives the rank
identity. Both $NP(L)$ and $N(1-P)(L)$ lie in $L$, so
$NL\subset G\oplus B$ and the index is finite.
The maps induced by $P$ and $1-P$ give isomorphisms
\[
 L/(G\oplus B)\simeq P(L)/G
                  \simeq (1-P)(L)/B.
\]
For example, $Px\in G$ implies $x-Px\in B$, which identifies the first
kernel. The middle quotient embeds in $N^{-1}G/G$ and the last in
$N^{-1}B/B$. As abelian groups, $G$ and $B$ are free of ranks $dg$
and $db$, so these latter quotients have orders $N^{dg}$ and $N^{db}$.
This proves the index bound.

Orthogonality at all infinite places gives
\[
 \operatorname{covol}(G)\operatorname{covol}(B)
       =I\operatorname{covol}(L).
\]
The discriminant terms in the definition of degree add with the ranks.
Taking logarithms therefore gives
$\widehat{\deg}L=\widehat{\deg}G+\widehat{\deg}B+d^{-1}\log I$,
which is the asserted slope identity. If $b=0$, then $G=L$ and $I=1$.
\end{proof}

Write $\nu(E)=2\mu(E)$, the slope associated with squared norms.
The factor two will be useful when logarithmic covolumes are expressed
in the same convention.

\begin{corollary}[A covolume bound from the discarded degree]
\label{cs-exceptional-budget}
Let $T,q$ be real, let $C\ge0$ and $U>0$, and suppose
\[
 q\le-\nu(G),\qquad \nu(L)\ge-T-CU.
\]
For $b>0$ define
\[
 \mathcal D=
 \frac bg\bigl(\nu(B)+T+CU\bigr)+\frac{2\log I}{dg};
\]
put $\mathcal D=0$ when $b=0$. Then
\[
 q\le T+CU+\mathcal D.
\]
In particular, $\mathcal D\le\kappa U$ implies
$q\le T+(C+\kappa)U$ for every $\kappa\ge0$.
If $\nu(B)\le\Lambda$, then, with $\eta=b/r<1$,
\[
 \mathcal D\le
 \frac{\eta}{1-\eta}\bigl(\Lambda+T+CU\bigr)
 +2\min\!\left\{\frac{\eta}{1-\eta},1\right\}\log N.
\]
For $b=0$ this last inequality holds without any condition on $\Lambda$.
\end{corollary}

\begin{proof}
Twice the identity in Proposition~\ref{cs-orthogonal-cut} gives
\[
 -\nu(G)=-\frac rg\nu(L)+\frac bg\nu(B)+\frac{2\log I}{dg}.
\]
Insert the lower bound for $\nu(L)$ and use $r/g=1+b/g$.
The remaining estimate follows from the index bound and
$b/g=\eta/(1-\eta)$.
\end{proof}

Only the mean slope of the discarded lattice enters this bound.
An upper bound for the maximal slope of a larger lattice containing
$L$ is one way to control it, but is stronger than required.
Likewise, no prescribed rate of decay of the exceptional rank occurs
in Corollary~\ref{cs-exceptional-budget}.

\begin{corollary}[Polynomial decay can suffice]\label{cs-polynomial-exceptional}
Suppose the data of Corollary~\ref{cs-exceptional-budget} vary with a
real parameter $x\to\infty$, possibly with further indices, and that
$\nu(B)\le\Lambda$ whenever $b>0$. Assume, uniformly in those indices,
that for fixed $a>p\ge0$ and constants $c,c'>0$,
\[
 \eta\le c x^{-a},\qquad
 \frac{|\Lambda|+|T|+CU+2\log N}{U}\le c'x^p.
\]
Then, uniformly in those indices, for all sufficiently large $x$,
\[
 \mathcal D/U\le 2cc'x^{p-a}.
\]
Consequently, for each $\kappa>0$, the bound
$q\le T+(C+\kappa)U$ holds for all sufficiently large $x$.
\end{corollary}

\begin{proof}
Eventually $\eta\le1/2$, and
$\eta/(1-\eta)\le2c x^{-a}$. The upper bound in
Corollary~\ref{cs-exceptional-budget}, divided by $U$, is therefore
at most $2cc'x^{p-a}$, which tends to zero.
\end{proof}

Small exceptional rank alone is insufficient. Over $\mathbb Q$, take
an orthogonal lattice with $n$ basis vectors of norm $e^n$ and one
of norm $e^{-n^2}$. Project onto the first $n$ coordinates. Here
$N=I=1$, $\mu(L)=0$ and $b/r=1/(n+1)\to0$, but
$\mu(G)=-n$. The discarded mean slope is $n^2$ and must be controlled.
The index term is also necessary: for
$L=\mathbb Ze_1+\mathbb Z(e_1+e_2)/N$ with the standard metric and
projection onto $\mathbb Qe_1$, both $G$ and $B$ have slope zero,
whereas $I=N$ and $\mu(L)=\tfrac12\log N$.

\subsection{Weighted lower estimates}

\begin{lemma}[Weighted lower estimates]\label{cs-weighted-lower}
Let $I$ be a finite nonempty set with positive weights, and write a bar
for its weighted mean. Let $\nu_i,\mu_i,h_i,D_i,E_i,\ell_i$ be real
numbers, and let $c,F,\Gamma,A,\mathcal M,P$ be real. Suppose
\[
 \overline{\nu-\mu+E+\ell}\ge\Gamma+F,
 \qquad c+\overline h-F\le\overline D,
 \qquad \overline{\nu-h}-c\le P.
\]
Then
\[
 A-\mathcal M\le A-\Gamma+
       \overline{D+E+\ell-\mu}-\mathcal M+P.
\]
\end{lemma}
\begin{proof}
Put $s=\overline{\nu-h}$. The first two assumptions give
\[
 s\ge\Gamma+F+\overline{\mu-E-\ell-h}
   \ge\Gamma+c-\overline{D+E+\ell-\mu}.
\]
Use $s-c\le P$ and rearrange.
\end{proof}

If these data are indexed in turn by a finite nonempty translation set,
one may average the resulting inequalities over that set. Thus only the
average of the final error expression must be controlled. The divisor
and lower quotient estimates themselves may be proved in weighted mean.
The quantity $P$ may be supplied directly by either slope comparison
above, with the same integral maps and common target lattice. Keeping
its actual value avoids a preliminary replacement of all frame defects
and height gaps by common upper bounds. In particular, any orthogonal
projection estimate used for its covolume term must concern that same
lattice, with the same metric and projection.

\subsection{Finite character averaging}

\begin{proposition}[Finite character averaging]\label{cs-finite-average}
Fix $x>0$. Let $A$ be a finite abelian group, $B$ a subgroup, $T\subset A$ and
$S\subset B$ nonempty finite sets. Write $n_T=|T|$, $d=|S|/|A|$, and
\[
 G=(\mathbf Z/2\mathbf Z)^2\times A,\quad
 \Phi=\{((0,e),a):e\in\mathbf Z/2\mathbf Z,\ a\in A\},\quad
 f=|\Phi|.
\]
Put $c=((1,1),0)$, $t_0=((0,0),t)$, $e_1(t)=((1,0),t)$ and
$\Phi_{S,u}=\{((0,1),s+u):s\in S\}$. Embed $B$ in $G$ by
$u\mapsto((0,0),u)$, and let
$X=B^\perp\subset\widehat G$.
Suppose $m\ge1$ is an integer, $b>0$, $a_t>0$ and $w_{u,i}(\rho)>0$ for
$u\in B$, $t\in T$, $1\le i\le m$ and $\rho\in\Phi$.
Let $Z_\psi:(0,\infty)\to\mathbb C$ be functions, and suppose real
numbers $E_{u,t,i}$ satisfy
\[
 E_{u,t,i}=\Re\frac1f\sum_{g\notin\Phi}\sum_{\psi\in\widehat G}
       \psi(-g)w_{u,i}(c+g+t_0)
          Z_\psi\!\left(\frac{a_t}{w_{u,i}(c+g+t_0)}\right).
\]
Choose $\mathcal K\subset X$ and positive weights $\lambda_\psi$, and put
\[
 H(x)=\sum_{\psi\in\mathcal K}\lambda_\psi\Re Z_\psi(x),\qquad
 \overline E=\frac1{|B|}\sum_{u\in B}\sum_{t\in T}\sum_{i=1}^m E_{u,t,i},
 \qquad
 \alpha_\psi=\frac1{n_T}\sum_{t\in T}\psi(e_1(t)).
\]
In the transfer sum, write $\rho=c+g+t_0$. Define
\begin{align*}
 R_F(x)={}&\sum_{\psi\in\mathcal K}|\alpha_\psi-\lambda_\psi|\,|Z_\psi(x)|
       +\sum_{\psi\in X\setminus\mathcal K}|\alpha_\psi|\,|Z_\psi(x)|\\
 &+\frac1{n_T}\sum_{t\in T}\sum_{\psi\in X}
             |Z_\psi(a_t/b)-Z_\psi(x)|,\\
 R_T={}&\frac1{|B|f}\sum_{u\in B}\sum_{t\in T}\sum_{i=1}^m
           \sum_{g\notin\Phi}\sum_{\psi\in\widehat G}
 \left|w_{u,i}(\rho)Z_\psi\!\left(\frac{a_t}{w_{u,i}(\rho)}\right)
   -\mathbf1_{\Phi_{S,u}}(\rho)bZ_\psi(a_t/b)\right|.
\end{align*}
Then
\[
                H(x)\le R_F(x)+\frac{2(\overline E+R_T)}{n_Tmbd}.
\]
One may replace $R_F$ by any upper bound for the actual real projection
error
\[
 H(x)-\frac1{n_T}\sum_{t\in T}
          \Re\sum_{\psi\in X}\psi(e_1(t))Z_\psi(a_t/b).
\]
\end{proposition}
\begin{proof}
The triangle inequality bounds the last displayed projection error by
$R_F$. In the defining expression for $E_{u,t,i}$, replacing each weighted
function by its thick-part value costs at most the corresponding summand
of $R_T$ after averaging. On the thick part,
$g=e_1(s+u-t)$ with $s\in S$.
Summation over $u\in B$ annihilates every character outside $X$; for
$\psi\in X$ the remaining phase is $\psi(e_1(t))$, since $s\in B$.
The factor from summing over $s$ and dividing by $f=2|A|$ is $d/2$.
Consequently
\[
 \frac{n_Tmbd}{2}\bigl(H(x)-R_F(x)\bigr)
        \le\overline E+R_T.
\]
The coefficient on the left is positive, so division proves the result.
The argument used only an upper bound for the real projection error,
which proves the final assertion.
\end{proof}

When the functions $Z_\psi$ are the smoothed Dirichlet sums defined in
Section~\ref{sec:purpose}, this proposition gives a sufficient arithmetic
budget for Corollary~\ref{cor:varying-positive-families}. For each fixed
accuracy, control the projection error and the normalized sum of the
averaged discrepancy and transfer error by powers with total exponent
loss tending to zero. The same selected positive combination must be
retained throughout that accuracy's sequence of whole intervals, even
if the ambient finite groups and auxiliary data grow with $x$.
The proposition imposes no separate square-root choice of the chain
length $m$; such a choice is one way to balance the errors in a particular
construction.

\subsection{A convenient stronger set of arithmetic estimates}

Here is a precise reduction of the remaining numerical comparison.
It avoids auxiliary parameters that disappear by cancellation.

\begin{lemma}[Combination of the local estimates]\label{cs-combination}
Let $I$ be a nonempty finite set, with positive weights $d_i$ and
weighted mean denoted by a bar. Let
$\nu_i\le\nu_i^+$, $h_i$, $M_i$, $M_i^+$, $D_i$, $E_i$, and $\ell_i$
be real numbers. The scalar parameters
$c,F,G,A,\mathcal M,K,M_-,M_+,E_+,V,L,L_0$ are also real.
Suppose
\begin{align*}
 c+h_i-F&\le D_i,&
 \nu_i&\ge M_i+G+F-E_i-\ell_i,\\
 \nu_i^+&\le M_i^++G+F-E_i+E_+,&
 -M_-&\le M_i,\qquad M_i^+\le M_+,\\
 D_i+E_i&\le\mathcal M+K,&
 (-E_i-h_i)-(-E_j-h_j)&\le V,\\
 A-G&\le L,&\overline\ell&\le L_0.
\end{align*}
Put $s=\overline{\nu-h}$ and
$U=\max_i(\nu_i^+-h_i)-s$. Suppose $H\ge0$ and
\[
 s-c\le\frac{2U\log(1+H/U)+(\eta+4\sqrt e)H}{1-\delta},
 \qquad 0\le\delta<1,\quad\eta,e\ge0.
\]
Then, with $B=M_++M_-+V+E_++L_0$, one has $0\le U\le B$ and
\[
 A-\mathcal M\le L+K+M_-+L_0+
 \frac{2B\log(1+H/B)+(\eta+4\sqrt e)H}{1-\delta}.
\]
If $H\le Xb$ with $X\ge0$, $b>0$, the last numerator is at most
$2(B+b)\log(1+X)+(\eta+4\sqrt e)Xb$.
\end{lemma}
\begin{proof}
The first, second, fifth and lower-slope inequalities give
$s\ge G+c-\mathcal M-K-M_--\overline\ell$.
Comparing the upper bound for $\nu_i^+-h_i$ with the weighted lower
bound for $\nu_j-h_j$ gives
$U\le M_++M_-+V+E_++\overline\ell\le B$.
Also $U\ge0$, since $\nu_i\le\nu_i^+$.
The function $z\mapsto z\log(1+H/z)$ is increasing for $z>0$:
its derivative is
$\log(1+H/z)-H/(H+z)\ge0$, by integrating
$1/(1+t)\ge1/(1+H/z)$ from $0$ to $H/z$.
The asserted inequality now follows by addition. Finally, concavity of
the logarithm, with weights $B/(B+b)$ and $b/(B+b)$, gives
\[
 B\log(1+H/B)\le(B+b)\log(1+H/(B+b))
 \le(B+b)\log(1+X).
\]
The limiting cases $B=0$ or $H=0$ follow by continuity.
\end{proof}

For comparison with the proposed arithmetic construction, let $m\ge1$ be
its integer number of lifts, let $\mathcal T$ be the finite nonempty
index set of interpolation orders $a_t>0$, and let $T_i>0$ be its
auxiliary primes, $g>0$ its product of two local factors, and $w_{j\rho}>0$
its weights, where $f\ge1$ is an integer,
$1\le j\le m$, and $1\le\rho\le f$. Write $\mathcal R^\vee$
for the target Hermitian lattice, $h_i$ for the map-height bounds,
$F\in\mathbb R$ for the local correction, and $v>0$ for the
archimedean normalization. Set
\[
 S=\sum_t a_t,\qquad
 \log\Omega=\gamma+\log(\pi/v)-\ell,
\]
where $\gamma$ is Euler's constant and $\ell$ is the relevant real
logarithmic derivative of a quadratic $L$-function at $1$. The required
divisor inequality has the concrete form
\begin{equation}\label{cs-required-divisor}
 \mu(\mathcal R^\vee)+h_i-F\le
 \sum_t\{m(a_t\log a_t-a_t-a_t\log\Omega)-a_t\log T_i\}
       +\mathcal E_D.
\end{equation}
The slopes here and below may all be multiplied by the same fixed
normalizing factor; the preceding homogeneous inequalities are unchanged.
Let $\mathcal L,B_0,N_c\ge0$ denote the error parameters whose
arithmetic interpretation and simultaneous bounds must be supplied by
that construction. The proposed bound for the error is
\[
 \mathcal E_D=C|\mathcal T|a\mathcal L
             +Cm(B_0+N_c)x^{2\varepsilon},
 \qquad a=xb,\quad b=\exp(x^\varepsilon),\quad
 x^{1/2-\varepsilon}\le m\le x^{1/2+\varepsilon/2}.
\]
The multiplicative constant $C$ may depend on the family and on
$\varepsilon$, while the exponent $2$ must be independent of
$\varepsilon$. The assertion is required for every word and all
sufficiently large parameters in the selected family.
In particular, a multiplicative constant depending on $\varepsilon$
can be absorbed by increasing that threshold. An exponent coefficient
growing like $1/\varepsilon$ cannot be treated in the same way.

The remaining cancellation is elementary. Suppose
\begin{align*}
 E_i&\le mS\log g+\mathcal E_E,\\
 \log T_i&\ge
   \sum_{j=1}^m\frac1f\sum_{\rho=1}^f\log(w_{j\rho}v/\pi)
       +m\log g-\mathcal E_C.
\end{align*}
If $D_i$ denotes the right side of \eqref{cs-required-divisor}, then
\[
 D_i+E_i\le\mathcal M+S\mathcal E_C+\mathcal E_D+\mathcal E_E,
\]
where
\[
 \mathcal M=\sum_{j=1}^m\sum_t
 a_t\left(\log a_t-1-\gamma+\ell
                   -\frac1f\sum_{\rho=1}^f\log w_{j\rho}\right).
\]
This follows by direct substitution: both $mS\log g$ and
$mS\log(v/\pi)$ cancel. Together with comparison of the source slopes,
the mean loss bound, and Proposition~\ref{cs-slope},
Lemma~\ref{cs-combination} proves the final scalar estimate. It removes
the term $\log|I|$ and replaces the coefficient $4$ of
$U\log(1+H/U)$ in the displayed source comparison by $2$.
These improvements simplify constants and exposition. They do not
prove \eqref{cs-required-divisor}, nor the matrix estimates for the
automorphic maps, nor by themselves the required bounds for the selected
twisted-character sums.  The latter involve a fixed finite family of
base characters; only the RH subcase reduces to three quadratic
characters.  A separate argument by averaging over the auxiliary primes
can establish this selection under a suitable uniform QRH hypothesis.

The source of \eqref{cs-required-divisor} is more than the
nonvanishing of one integer determinant. At each finite place it needs
pairing divisibility for every vector in a prescribed lattice, and a
bound for a rank-normalized sum of valuations of bordered minors.
At every infinite place it needs upper bounds for operator norms of a
chain of rectangular matrices, together with lower bounds for the
eigenvalues in the quotient metrics. These are separate assertions
about the constructed matrices. The elementary determinant of the
preceding quadratic argument does not establish them.

\begin{remark}[A counterexample to a possible simplification]\label{cs-counterexample}
Nonvanishing modulo every prime does not imply the matrix estimate
\eqref{cs-isotropic} with a small $\eta$. For any integer $M\ge1$, take
the columns of the unimodular matrix
\[
 \begin{pmatrix}1&M\\0&1\end{pmatrix}
\]
and normalize each column to length $1$. Thus
$A_1=(1,0)^t$ and $A_2=(M,1)^t/\sqrt{M^2+1}$.
Both diagonal errors vanish, but
\[
 \operatorname{tr}Q=2,\qquad
 \operatorname{tr}(Q^2)=2+\frac{2M^2}{M^2+1},\qquad
 \frac{2\operatorname{tr}(Q^2)}{(\operatorname{tr}Q)^2}-1
       =\frac{M^2}{M^2+1}.
\]
The smallest permissible $\eta$ tends to $1$. This is a counterexample
to replacing the theta estimates by determinant nonvanishing; it is
not a counterexample to the theta estimates asserted for the specific
automorphic construction.
\end{remark}


